\documentclass[a4paper]{article}
\usepackage{todonotes}

\title{Optical Fibres and Optocommunication Laboratory 5:
Erbium Doped Fiber Amplifier (EDFA) --- parameters and characteristics
}
\author{Kamil Czop 259613\\Sergiusz Warga 230757}
\date{\today}

\begin{document}

\maketitle

\section{Theoretical introduction}

Amplified Spontaneous Emission

\todo[inline]{Describe the construction of the EDFA}

The WDM coupler ensures no losses, as it combines the signals of different wavelengths.
The WDM coupler allows us to combine two signals of two wavelengths with almost no losses.

In our laboratory setup we have forward-propagating pump and backward-propagating pump. 

Seed laser has the Power of 5 dBm.

\todo[inline]{Describe the principle of operation of the EDFA}

The emission band of the EDFA coincides with the third telecommunication window.

\todo[inline]{Whe the fibre is emitting a green light?}




\section{Experiments}

\subsection{Seting up}
\subsection{Power measurements}
Reference level.
\subsection{Amplifier characteristics}
\subsection{Noise Figure}
Analyser resolution 0.1 nm 

lambda = 1554.520 nm

We want to perform the measurements for 0 dBm and 40 dBm of the input power.




\end{document}
